% Toggles between whether printed as a handout or "bullet-per-slide"
\documentclass{beamer}
% \documentclass[handout]{beamer}
\usepackage[utf8]{inputenc}
\usetheme{Warsaw}
\usecolortheme{uiowa}
\useoutertheme[subsection=false]{smoothbars}

\usepackage{array}
\usepackage{booktabs}
\usepackage{graphics}
\usepackage{natbib}
\usepackage{xcolor}
\usepackage{amsmath}
\usepackage{amssymb}
\usepackage[english]{babel}
\usepackage{caption}
\usepackage{amsfonts}
\usepackage{tikz}
\usepackage{bm}
\usepackage{hyperref}
\newcommand\blfootnote[1]{%
  \begingroup
  \renewcommand\thefootnote{}\footnote{#1}%
  \addtocounter{footnote}{-1}%
  \endgroup
}

% Uppercase roman letters
\providecommand{\A}{\mathbf{A}}
\providecommand{\B}{\mathbf{B}}
\providecommand{\C}{\mathbf{C}}
\providecommand{\D}{\mathbf{D}}
\providecommand{\E}{\mathbf{E}}
\providecommand{\F}{\mathbf{F}}
\providecommand{\G}{\mathbf{G}}
\renewcommand{\H}{\mathbf{H}}
\providecommand{\I}{\mathbf{I}}
\providecommand{\J}{\mathbf{J}}
\providecommand{\K}{\mathbf{K}}
\providecommand{\L}{\mathbf{L}}
\providecommand{\M}{\mathbf{M}}
\providecommand{\N}{\mathbf{N}}
\providecommand{\O}{\mathbf{O}}
\renewcommand{\P}{\mathbf{P}}
\providecommand{\Q}{\mathbf{Q}}
\providecommand{\R}{\mathbf{R}}
\renewcommand{\S}{\mathbf{S}}
\providecommand{\T}{\mathbf{T}}
\providecommand{\U}{\mathbf{U}}
\providecommand{\V}{\mathbf{V}}
\providecommand{\W}{\mathbf{W}}
\providecommand{\X}{\mathbf{X}}
\providecommand{\Y}{\mathbf{Y}}
\providecommand{\Z}{\mathbf{Z}}

% Lowercase roman letters (c, d, u, v have to be treated special; see end)
\renewcommand{\a}{\mathbf{a}}
\renewcommand{\b}{\mathbf{b}}
\newcommand{\e}{\mathbf{e}}
\newcommand{\f}{\mathbf{f}}
\newcommand{\g}{\mathbf{g}}
\newcommand{\h}{\mathbf{h}}
\renewcommand{\i}{\mathbf{i}}
\renewcommand{\j}{\mathbf{j}}
\renewcommand{\k}{\mathbf{k}}
\renewcommand{\l}{\mathbf{l}}
\newcommand{\m}{\mathbf{m}}
\newcommand{\n}{\mathbf{n}}
\renewcommand{\o}{\mathbf{o}}
\newcommand{\p}{\mathbf{p}}
\newcommand{\q}{\mathbf{q}}
\renewcommand{\r}{\mathbf{r}}
\newcommand{\s}{\mathbf{s}}
\renewcommand{\t}{\mathbf{t}}
\newcommand{\w}{\mathbf{w}}
\newcommand{\x}{\mathbf{x}}
\newcommand{\y}{\mathbf{y}}
\newcommand{\z}{\mathbf{z}}
\let\origc\c
\DeclareRobustCommand\c{\ifmmode\mathbf{c}\else\expandafter\origc\fi}
\let\origd\d
\DeclareRobustCommand\d{\ifmmode\mathbf{d}\else\expandafter\origd\fi}
\let\origu\u
\DeclareRobustCommand\u{\ifmmode\mathbf{u}\else\expandafter\origu\fi}
\let\origd\v
\DeclareRobustCommand\v{\ifmmode\mathbf{v}\else\expandafter\origv\fi}

% Roman letters with hats
\providecommand{\fhat}{\hat{f}}
\providecommand{\shat}{\hat{s}}
\providecommand{\yhat}{\hat{y}}
\providecommand{\bfhat}{\hat{\bf f}}
\providecommand{\bshat}{\hat{\mathbf{s}}}
\providecommand{\byhat}{\hat{\mathbf{y}}}

% Roman letters with subscripts
\providecommand{\rj}{\mathbf{r}_{-j}}
\providecommand{\xj}{\mathbf{x}_{-j}}
\providecommand{\Xj}{\mathbf{X}_{-j}}

% Roman letters with tildes
\providecommand{\rr}{\tilde{r}}
\providecommand{\xx}{\tilde{x}}
\providecommand{\yy}{\tilde{y}}
\providecommand{\zz}{\tilde{z}}
\providecommand{\brr}{\tilde{\mathbf{r}}}
\providecommand{\bxx}{\tilde{\mathbf{x}}}
\providecommand{\byy}{\tilde{\mathbf{y}}}
\providecommand{\bzz}{\tilde{\mathbf{z}}}
\providecommand{\XX}{\widetilde{\mathbf{X}}}
\providecommand{\WW}{\widetilde{\mathbf{W}}}
\providecommand{\ZZ}{\tilde{\mathbf{Z}}}

% Script letters
\providecommand{\cA}{\mathcal{A}}
\providecommand{\cB}{\mathcal{B}}
\providecommand{\cC}{\mathcal{C}}
\providecommand{\cD}{\mathcal{D}}
\providecommand{\cE}{\mathcal{E}}
\providecommand{\cF}{\mathcal{F}}
\providecommand{\cG}{\mathcal{G}}
\providecommand{\cH}{\mathcal{H}}
\providecommand{\cI}{\mathcal{I}}
\providecommand{\cM}{\mathcal{M}}
\providecommand{\cN}{\mathcal{N}}
\providecommand{\cR}{\mathcal{R}}
\providecommand{\cS}{\mathcal{S}}
\providecommand{\cT}{\mathcal{T}}
\providecommand{\cU}{\mathcal{U}}
\providecommand{\cV}{\mathcal{V}}
\providecommand{\cW}{\mathcal{W}}
\providecommand{\cX}{\mathcal{X}}
\providecommand{\cY}{\mathcal{Y}}
\providecommand{\cZ}{\mathcal{Z}}

% Greek letters
\providecommand{\ah}{\hat{\alpha}}
\providecommand{\ba}{\boldsymbol{\alpha}}
\providecommand{\bah}{\hat{\boldsymbol{\alpha}}}
\providecommand{\bb}{\boldsymbol{\beta}}
\providecommand{\bh}{\widehat{\beta}}
\providecommand{\bbh}{\widehat{\boldsymbol{\beta}}}
\providecommand{\bd}{\boldsymbol{\delta}}
\renewcommand{\dh}{\hat{\delta}}
\providecommand{\be}{\boldsymbol{\eta}}
\providecommand{\eh}{\hat{\eta}}
\providecommand{\beh}{\hat{\boldsymbol{\eta}}}
\providecommand{\eps}{\epsilon}
\providecommand{\veps}{\varepsilon}
\providecommand{\bep}{\boldsymbol{\epsilon}}
\providecommand{\bvep}{\boldsymbol{\varepsilon}}
\providecommand{\eph}{\hat{\epsilon}}
\providecommand{\beph}{\hat{\boldsymbol{\epsilon}}}
\providecommand{\gam}{\gamma}
\providecommand{\bg}{\boldsymbol{\gamma}}
\providecommand{\bgh}{\hat{\boldsymbol{\gamma}}}
\providecommand{\bG}{\boldsymbol{\Gamma}}
\providecommand{\gh}{\hat{\gamma}}
\providecommand{\bm}{\boldsymbol{\mu}}
\providecommand{\mh}{\hat{\mu}}
\providecommand{\bmh}{\hat{\boldsymbol{\mu}}}
\providecommand{\lam}{\lambda}
\providecommand{\Lam}{\Lambda}
\providecommand{\bl}{\boldsymbol{\lambda}}
\providecommand{\lh}{\hat{\lambda}}
\providecommand{\bL}{\boldsymbol{\Lambda}}
\providecommand{\bO}{\boldsymbol{\Omega}}
\providecommand{\bp}{\boldsymbol{\pi}}
\providecommand{\ph}{\hat{\pi}}
\providecommand{\bph}{\hat{\boldsymbol{\pi}}}
\providecommand{\bpsi}{\boldsymbol{\psi}}
\providecommand{\sh}{\hat{\sigma}}
\providecommand{\bS}{\boldsymbol{\Sigma}}
\providecommand{\bSh}{\hat{\boldsymbol{\Sigma}}}
\providecommand{\bt}{\boldsymbol{\theta}}
\renewcommand{\th}{\hat{\theta}}
\providecommand{\ts}{\theta^*}
\providecommand{\bth}{\hat{\boldsymbol{\theta}}}
\providecommand{\bts}{\boldsymbol{\theta}^*}
\providecommand{\bT}{\boldsymbol{\Theta}}
\providecommand{\btau}{\boldsymbol{\tau}}
\providecommand{\tauh}{\hat{\tau}}
\providecommand{\btauh}{\hat{\boldsymbol{\tau}}}
\providecommand{\bbj}{\boldsymbol{\beta}_{-j}}
\providecommand{\bbhj}{\widehat{\boldsymbol{\beta}}_{-j}}

%% Greek letters with tildes
\providecommand{\betat}{\widetilde{\beta}}
\providecommand{\bbjt}{\widetilde{\boldsymbol{\beta}}_{-j}}
\providecommand{\bbt}{\widetilde{\boldsymbol{\beta}}}
\providecommand{\etat}{\tilde{\eta}}
\providecommand{\bet}{\widetilde{\boldsymbol{\eta}}}
\providecommand{\lamt}{\tilde{\lambda}}
\providecommand{\mt}{\tilde{\mu}}
\providecommand{\bmt}{\tilde{\boldsymbol{\mu}}}
\providecommand{\nut}{\tilde{\nu}}
\providecommand{\thetat}{\tilde{\theta}}
\providecommand{\btt}{\tilde{\boldsymbol{\theta}}}

% special rotated character symbols
\providecommand{\tone}{\widetilde{1}}
\providecommand{\bone}{{\mathbf{1}}}
\providecommand{\tbone}{\widetilde{\mathbf{1}}}
\providecommand{\tvep}{\widetilde{\varepsilon}}
\providecommand{\tveph}{\hat{\widetilde{\varepsilon}}}
\providecommand{\tbvep}{\widetilde{\boldsymbol{\varepsilon}}}
\providecommand{\ty}{\widetilde{y}}
\providecommand{\tx}{\widetilde{x}}

%% Operators
\renewcommand{\Pr}{\mathbb{P}}
\providecommand{\pr}{\textrm{Pr}}
\providecommand{\Ex}{\mathbb{E}}
\providecommand{\ex}{\textrm{E}}
\providecommand{\exh}{\widehat{\textrm{E}}}
\providecommand{\Var}{\mathbb{V}}
\providecommand{\var}{\textrm{Var}}
\providecommand{\varh}{\widehat{\textrm{Var}}}
\providecommand{\cov}{\textrm{Cov}}
\providecommand{\cor}{\textrm{Cor}}
\providecommand{\trace}{\textrm{tr}}
\providecommand{\rank}{\textrm{rank}}
\providecommand{\Norm}{\textrm{N}}
\providecommand{\Gmma}{\textrm{Gamma}}
\providecommand{\Beta}{\textrm{Beta}}
\providecommand{\Pois}{\textrm{Pois}}
\providecommand{\NBin}{\textrm{NegBin}}
\providecommand{\Unif}{\textrm{Unif}}
\providecommand{\Binom}{\textrm{Binom}}
\providecommand{\Bern}{\textrm{Bern}}
\providecommand{\Weib}{\textrm{Weibull}}
\providecommand{\Scx}{\textrm{Scaled-}\chi^2}
\providecommand{\Wish}{\textrm{Wishart}}
\providecommand{\Multinom}{\textrm{Multinom}}
\providecommand{\Dirch}{\textrm{Dir}}
\providecommand{\Exp}{\textrm{Exp}}
\providecommand{\diag}{\textrm{diag}}
\providecommand{\Tr}{^{\scriptscriptstyle\top}}

%% Statistical
\renewcommand{\l}{\ell}
\providecommand{\inP}{\overset{\textrm{P}}{\longrightarrow}}
\providecommand{\inAS}{\overset{\textrm{as\,}}{\longrightarrow}}
\providecommand{\inD}{\overset{\textrm{d}}{\longrightarrow}}
\providecommand{\inQM}{\overset{\textrm{qm\,}}{\longrightarrow}}
\providecommand{\adist}{\overset{\text{.}}{\sim}}
\providecommand{\AIC}{\textrm{AIC}}
\providecommand{\BIC}{\textrm{BIC}}
\providecommand{\CI}{\textrm{CI}}
\providecommand{\CV}{\textrm{CV}}
\providecommand{\FWER}{\textrm{FWER}}
\providecommand{\FDR}{\textrm{FDR}}
\providecommand{\FDRh}{\widehat{\textrm{FDR}}}
\providecommand{\GCV}{\textrm{GCV}}
\providecommand{\KL}{\textrm{KL}}
\providecommand{\ME}{\textrm{ME}}
\providecommand{\MLE}{\textrm{MLE}}
\providecommand{\MSE}{\textrm{MSE}}
\providecommand{\MSPE}{\textrm{MSPE}}
\providecommand{\PE}{\textrm{PE}}
\providecommand{\RSS}{\textrm{RSS}}
\providecommand{\SNR}{\textrm{SNR}}
\providecommand{\SE}{\textrm{SE}}
\providecommand{\SD}{\textrm{SD}}
\providecommand{\OR}{\textrm{OR}}
\providecommand{\RR}{\textrm{RR}}
\providecommand{\df}{\textrm{df}}
\providecommand{\odds}{\textrm{odds}}
\providecommand{\loglik}{\textrm{loglik}}
\providecommand{\logit}{\textrm{logit}}
\providecommand{\const}{\textrm{constant}}
\providecommand{\OLS}{\textrm{OLS}}

% Fisher/observed information
\providecommand{\oI}{\boldsymbol{\mathcal{I}}}
\providecommand{\oV}{\boldsymbol{\mathcal{V}}}
\providecommand{\fI}{\mathbscr{I}}
\providecommand{\fV}{\mathbscr{V}}

% Independence
\providecommand{\iid}{\overset{\text{iid}}{\sim}}
\providecommand{\ind}{\perp\!\!\!\perp}
\providecommand{\IID}{\overset{\scriptscriptstyle \ind}{\sim}}

%% Mathematical
\providecommand{\0}{\mathbf{1}}
\providecommand{\1}{\mathbf{0}}
\providecommand{\one}{\mathbf{1}}
\providecommand{\zero}{\mathbf{0}}
\providecommand{\intii}{\int_{-\infty}^{\infty}}
\providecommand{\real}{\mathbb{R}}
%%\providecommand{\Real}{\mathds{R}} %% Requires dsfonts
\providecommand{\dl}{d_-}
\providecommand{\dr}{d_+}
\providecommand{\sign}{\textrm{sign}}
\providecommand{\given}{\,\vert\,}
\providecommand{\abs}[1]{\left\lvert#1\right\rvert}
\providecommand{\sabs}[1]{\lvert#1\rvert}
\providecommand{\norm}[1]{\lVert#1\rVert}
\providecommand{\bignorm}[1]{\left\lVert#1\right\rVert}
\providecommand{\gl}{\prec}
\providecommand{\gle}{\preceq}
\providecommand{\gg}{\succ}
\providecommand{\gge}{\succeq}
\DeclareMathOperator*{\argmin}{\arg\min}
\DeclareMathOperator*{\argmax}{\arg\max}
\providecommand{\pf}[2]{\frac{\partial #1}{\partial #2}}

%% Equations
\providecommand{\al}[2]{\begin{align}\label{#1}#2\end{align}}
\providecommand{\as}[1]{\begin{align*}#1\end{align*}}
\providecommand{\als}[2]{\begin{align}\label{#1}\begin{split}#2\end{split}\end{align}}

% misc other commands ----------------

\def\be{\begin{equation*}}
\def\ee{\end{equation*}}

\def\E{{\mbox E}}

\def\thhat{\hat{\theta}}


% link colors
\hypersetup{
  colorlinks=true,
  citecolor=black,
  linkcolor=.,
  urlcolor=uiblue
}

% ------------------------------------------------------------------
% Title info
% ------------------------------------------------------------------
\title[Collaboration with GitHub \hspace{3mm} Version Control in Practice \hspace{7em} \insertframenumber/\inserttotalframenumber]
{Version Control, GitHub Collaboration,\\ and Reproducible Team Workflow}

\author[Group Project]{Dylan Day, Shuo Li, Trevor Debutch, Shouhardyo Sarkar}
\date{\today}
\institute[The University of Iowa]{The University of Iowa, Department of Statistics and Actuarial Science}

\logo{\includegraphics[width=0.12\textwidth]{iowa-logo.png}}

% ------------------------------------------------------------------
\begin{document}
% ------------------------------------------------------------------

\frame{\titlepage}

% ==============================
\section{Motivation \& Goals}
% ==============================

\begin{frame}{Why Are We Talking About Version Control?}
  \begin{itemize}
    \item We work in groups and need to edit the same files without overwriting each other.
    \item We want a record of \textbf{who changed what and when}.
    \item We want to be able to \textbf{go back in time} if something breaks.
    \item We want our work to be \textbf{reproducible} on any machine.
  \end{itemize}

  \vspace{2mm}
  \textbf{Version control} is the system that tracks changes to files over time so you can recall or restore specific versions later.
\end{frame}

\begin{frame}{What We'll Show in This Talk}
  \begin{itemize}
    \item What Git and GitHub actually do (not just buzzwords).
    \item How teams collaborate using branches, commits, and pull requests.
    \item How to connect a local folder to a GitHub repo and push/pull.
    \item How this supports small coding projects (for class or research).
  \end{itemize}
\end{frame}

% ==============================
\section{Git, GitHub, and Key Concepts}
% ==============================

\begin{frame}{Git vs GitHub}
  \begin{block}{Git}
    A distributed version control system: every person has a full copy of the repository, history and all.
    Fast, offline-friendly.
  \end{block}

  \begin{block}{GitHub}
    A cloud platform \textbf{built on top of Git} that adds:
    \begin{itemize}
      \item remote hosting / backup
      \item pull requests \& code review
      \item issues \& project boards
      \item CI/CD automation
    \end{itemize}
    GitHub turns source code into a collaborative workspace.
  \end{block}
\end{frame}

\begin{frame}{GitHub Glossary}
  \begin{itemize}
    \item \textbf{Repository (repo)}: the project folder on GitHub.
    \item \textbf{Commit}: a saved change with a message.
    \item \textbf{Branch}: a parallel version of the project used to develop features safely.
    \item \textbf{Pull Request (PR)}: a request to merge changes from one branch into another after review.
  \end{itemize}

  \vspace{2mm}
  These concepts let multiple people experiment, review, and merge without breaking \texttt{main}.
\end{frame}

\begin{frame}{Why Use GitHub in a Class Project?}
  \begin{itemize}
    \item \textbf{Collaboration}: everyone can contribute code, documentation, or data without emailing files.
    \item \textbf{History}: you can see exactly how the project evolved.
    \item \textbf{Accountability}: each commit is tied to an author.
    \item \textbf{Safety}: if someone breaks something, you can roll back.
    \item \textbf{Publishability}: you can keep it private or make it public to show future advisors / employers.
  \end{itemize}
\end{frame}

% ==============================
\section{Basic Git Workflow}
% ==============================

\begin{frame}{Core Commands You Will Actually Use}
  \begin{itemize}
    \item \texttt{git clone <repo>} \\
      download a copy of a GitHub repo to your computer
    \item \texttt{git add .} \\
      stage your edits (tell Git which changes you want to commit)
    \item \texttt{git commit -m "message"} \\
      record a snapshot of those changes
    \item \texttt{git pull} and \texttt{git push} \\
      sync with GitHub (pull down others' work, push up yours)
    \item \texttt{git branch}, \texttt{git merge} \\
      work on new ideas in isolation, then merge them back
  \end{itemize}

  \vspace{2mm}
  These actions are the backbone of team development.
\end{frame}

\begin{frame}{Branches and Pull Requests}
  \begin{itemize}
    \item Each teammate can create a new branch (for example, \texttt{feature-ui}).
    \item You do your work there without touching the \texttt{main} branch.
    \item When ready, you open a \textbf{Pull Request} on GitHub:
      \begin{itemize}
        \item show what changed
        \item explain why
        \item request review
      \end{itemize}
    \item After review, the branch is merged into \texttt{main}.
  \end{itemize}

  \vspace{2mm}
  This workflow gives you feedback, quality control, and a clean history instead of chaos.
\end{frame}

% ==============================
\section{Connecting Local and Remote}
% ==============================

\begin{frame}{Step 1: Clone the Repository}
  \begin{itemize}
    \item One person creates a repo on GitHub.
    \item Everyone else runs:
  \end{itemize}

  \vspace{1mm}
  \texttt{git clone https://github.com/owner/RepoName.git}

  \vspace{2mm}
  Now you have a full local copy of the project folder, including history and branches.
\end{frame}

\begin{frame}[fragile]{Step 2: Check the Remote Link}
  \begin{itemize}
    \item Inside the project folder, verify GitHub is set as ``origin'':
  \end{itemize}

  \vspace{1mm}
  \texttt{git remote -v}

  \vspace{2mm}
  Expected:
  \begin{verbatim}
origin  https://github.com/owner/RepoName.git (fetch)
origin  https://github.com/owner/RepoName.git (push)
  \end{verbatim}

  \vspace{2mm}
  If you see that, your local repo is correctly connected to the GitHub repo.
\end{frame}

\begin{frame}{Step 3: Pull and Push}
  \begin{itemize}
    \item \textbf{Pull} before you start editing:
  \end{itemize}

  \vspace{1mm}
  \texttt{git pull origin main}

  \vspace{3mm}
  \begin{itemize}
    \item After editing files:
  \end{itemize}

  \vspace{1mm}
  \texttt{git add .}\\
  \texttt{git commit -m "explain what you changed"}\\
  \texttt{git push origin main}

  \vspace{2mm}
  If push is rejected, it means someone else updated \texttt{main}; run \texttt{git pull} to merge first.
\end{frame}

\begin{frame}{Step 4: Authentication / Permissions}
  \begin{itemize}
    \item GitHub requires you to prove who you are when pushing.
    \item If you see ``Permission denied'' or ``Authentication failed'':
      \begin{itemize}
        \item You may need a Personal Access Token (PAT) instead of a password.
        \item Or you may not have been added as a collaborator yet.
      \end{itemize}
    \item Good news: pulling (reading) is often easier than pushing (writing).
  \end{itemize}
\end{frame}

% ==============================
\section{Team Project Angle}
% ==============================

\begin{frame}{How This Supports a Team Project}
  \begin{itemize}
    \item Everyone can contribute in parallel:
      \begin{itemize}
        \item Code / analysis scripts
        \item README documentation
        \item Example data or test cases
      \end{itemize}
    \item Changes are tracked and credited to specific authors.
    \item The repo becomes a ``single source of truth'' for the project.
    \item You can show the final repo (or a link) as part of your assignment deliverable.
  \end{itemize}
\end{frame}

\begin{frame}{What We Practiced as a Group}
  \begin{itemize}
    \item Setting up a shared GitHub repository.
    \item Cloning to each teammate's laptop.
    \item Making edits on branches and creating pull requests for review.
    \item Resolving merge conflicts.
    \item Documenting everything in \texttt{README.md} so it is reproducible later.
  \end{itemize}
\end{frame}

% ==============================
\section{Takeaways}
% ==============================

\begin{frame}{Key Takeaways}
  \begin{itemize}
    \item Version control = a time machine for your project.
    \item GitHub = collaboration, accountability, and backup.
    \item Branches + pull requests = safe teamwork without stepping on each other's work.
    \item Pull before you push.
    \item Clear commit messages and good README files make your work reusable.
  \end{itemize}
\end{frame}

\begin{frame}{Questions?}
  \centering
  \Large Thank you!
\end{frame}

% ------------------------------------------------------------------
\bibliographystyle{plainnat}
\bibliography{refs}

\end{document}
